Como
\[
T \propto \frac{1}{\sqrt{g}},
\]
o periodo em outro planeta escala por
\[
\frac{T_{\text{planeta}}}{T_{\text{Terra}}} = \sqrt{\frac{g_{\text{Terra}}}{g_{\text{planeta}}}}.
\]
Usando $g_{\text{Terra}}\approx 9.81\,\text{m/s}^2$, $g_{\text{Marte}}\approx 3.71\,\text{m/s}^2$ e $g_{\text{Plutao}}\approx 0.62\,\text{m/s}^2$:
\[
\frac{T_{\text{Marte}}}{T_{\text{Terra}}} \approx 1.626,
\qquad
\frac{T_{\text{Plutao}}}{T_{\text{Terra}}} \approx 3.978.
\]
Logo, o pendulo oscilaria cerca de $62.6\%$ mais lentamente em Marte e quase $4$ vezes mais lentamente em Plutao.

